% =======================================================
% 4. INTERFACCIA API E COMUNICAZIONE
% =======================================================
\section{Interfaccia API e Comunicazione}

Al fine di garantire il disaccoppiamento tra l'editor Markdown (Frontend) e l'elaboratore di prompt (Backend), la comunicazione avviene tramite [specificare il paradigma scelto, es. REST, RPC, GraphQL].

\subsection{Struttura dei Payload AI}
La comunicazione relativa alle operazioni di Intelligenza Artificiale (Riassunto, Traduzione, Riscrittura, Sei Cappelli, Distant Writing) richiede una struttura dati definita.
[Descrivere in dettaglio le interfacce o i DTO (Data Transfer Objects) utilizzati per la comunicazione client-server. Es. struttura della richiesta contenente operazione, testo contesto e prompt; struttura della risposta contenente esito e testo generato].

\subsection{Gestione degli Errori e Fallback}
[Descrivere come vengono gestite, normalizzate e comunicate all'utente le anomalie di rete, i timeout dovuti all'LLM e le risposte malformate].